\label{sec:experiments}

\subsection{Experimental Validation Criteria}

To validate CrossPoll's effectiveness, we establish the following criteria:

\begin{enumerate}
    \item \textbf{Primary criterion}: CrossPoll achieves $\geq 15\%$ higher mAP@50 than COCO transfer at median label budget ($B=100$)
    \item \textbf{Secondary criterion}: Gains persist across all label budgets ($B \in \{25, 50, 100, 200, 500\}$)
    \item \textbf{Robustness criterion}: Performance is stable across three random seeds (std. dev. $< 3\%$ mAP@50)
    \item \textbf{Scalability criterion}: Multi-source transfer outperforms single-crop transfer at all budgets
\end{enumerate}

\subsection{Hyperparameter Sensitivity Analysis}

While we adopt standard Ultralytics defaults, we acknowledge potential sensitivity to:

\begin{itemize}
    \item \textbf{Learning rate decay}: Fine-tuning with lr=0.0001 may be suboptimal for very small budgets. Planned sensitivity study: compare lr $\in \{0.00001, 0.0001, 0.001\}$ on $B=50$
    \item \textbf{Training epochs}: 150-epoch limit may underutilize large budgets. Planned extension: run epochs $\in \{50, 100, 150, 200\}$ for $B=500$
    \item \textbf{Augmentation intensity}: Joint training uses standard Ultralytics augmentation. Could experiment with aggressive/conservative augmentation settings
\end{itemize}

These analyses are deferred to Section~\ref{sec:future} due to computational constraints.

\subsection{Failure Mode Analysis}

We identify potential failure modes and mitigation strategies:

\begin{table}[H]
\centering
\caption{Anticipated Failure Modes and Mitigations}
\label{tab:failures}
\begin{tabular}{lll}
\toprule
\textbf{Failure Mode} & \textbf{Symptom} & \textbf{Mitigation} \\
\midrule
Catastrophic forgetting & CrossPoll mAP drops below Scratch & Lower fine-tune LR, reduce epochs \\
Overfitting at $B=25$ & High train, low val mAP & Increase augmentation, reduce model capacity \\
Domain shift too large & All methods fail at high budgets & Verify Kiwi flower distinctiveness \\
Seed variance too high & Std. dev. $> 5\%$ across seeds & Increase training data or seeds \\
\bottomrule
\end{tabular}
\end{table}

\subsection{Computational Cost and Efficiency}

Training 60 runs (4 methods $\times$ 5 budgets $\times$ 3 seeds) on single M1 Pro requires approximately 14--15 hours wall-clock time (sequential execution). Optimizations applied:

\begin{table}[H]
\centering
\caption{Training Optimizations Applied}
\label{tab:optimizations}
\begin{tabular}{lcc}
\toprule
\textbf{Optimization} & \textbf{Expected Speedup} & \textbf{Applied} \\
\midrule
cache=`ram' (vs disk I/O) & 20--40\% & \checkmark \\
cos\_lr=True (convergence) & 0\% direct, enables 150 epoch cap & \checkmark \\
workers=4 (optimal for macOS) & 5--10\% (vs workers=8) & \checkmark \\
deterministic=False & 5--10\% & \checkmark \\
epochs: 300 $\to$ 150 & 50\% & \checkmark \\
\midrule
\textbf{Combined Effect} & \textbf{$\approx$ 2--3×} & \\
\bottomrule
\end{tabular}
